\subsection{Universal covering of a loop-disentanglable space}

Let B be a topological space and let b be a point of B. Recall that \(\Lambda_{b}(\mathrm{B})\) denotes the subspace of \(\mathcal{C}_{c}(\mathbf{I};\mathrm{B})\) formed by the paths with origin b, equipped with the quotient topology of the compact-convergence topology. We denote by \(e_{\mathrm{B}}\colon\Lambda_{b}(\mathrm{B})\to\mathrm{B}\) the map which, to a path, associates its endpoint; it is continuous. Write \(\lambda_{b}(\mathrm{B})\) for the quotient space of \(\Lambda_{b}(\mathrm{B})\) by the strict homotopy relation and equip it with the quotient topology. Since two strictly homotopic paths have the same endpoint, the map \(e_{B}\) defines, by passing to the quotient, a continuous map \(\varepsilon_{\mathrm{B}}\colon\lambda_{b}(\mathrm{B})\to\mathrm{B}\) , called the endpoint map.

THEOREM 1. --- Let B be a connected and locally path-connected topological space and let b be a point of B. The following properties are equivalent:

\begin{enumerate}
\def\labelenumi{(\roman{enumi})}
\item
  The space B is loop-disentanglable.
\item
  There exists a non-empty covering of B which is simply path-connected.
\item
  The space \(\lambda_{b}(\mathrm{B})\) , equipped with the endpoint map \(\varepsilon_{\mathrm{B}}\colon\lambda_{b}(\mathrm{B})\to\mathrm{B}\) , is a covering of B.
\item
  The group \(\pi_1(B, b)\) is discrete for the admissible topology.
\item
  The group \(\pi_{1}(B,b)\) is discrete for the quotient topology of the compact convergence topology.
\end{enumerate}

Moreover, when these conditions are satisfied, the space \(\lambda_{b}(\mathrm{B})\) is simply path-connected, simply connected, Galois with group \(\pi_{1}(\mathrm{B},b)^{\circ}\) , and the pointed covering \((\lambda_{b}(\mathrm{B}),\varepsilon_{b})\) is a universal covering of the pointed space \((\mathrm{B},b)\) .

It follows from the definition of a loop-disentanglable space that the subgroup of \(\pi_{1}(B,b)\) reduced to the identity element is admissible if and only if B is loop-disentanglable. This proves the equivalence (i)⇔(iv). Moreover, the equivalence of properties (iv) and (v) follows from remarks 4 and 5 of III, p. 320.

(iv)⇒(ii). By III, p. 316, prop. 5, there exists a covering (E, p) of B and a point \(x \in E_{b}\) such that \(p_{*}(\pi_{1}(E, x)) = \{e\}\) ; in particular, \(\pi_{1}(E, x)\) is reduced to the identity element. The arc-connected component of x in E is then a nonempty covering of B (I, p. 80, cor. 1 of prop. 6), simply path-connected.

(ii)⇒(iii). Let E be a nonempty simply path-connected covering of B. Let x be a point of \(E_{b}\) (such exist because B is connected, I, p. 74, prop. 4). The projection \(q: E \to B\) induces a homeomorphism \(\Lambda_{x}(E) \to \Lambda_{b}(B)\) (III, p. 302, cor. 2 of prop. 3), whence, by passing to arc-connected components, a homeomorphism \(q_{*}: \lambda_{x}(E) \simeq \lambda_{b}(B)\) . The endpoint map \(\Lambda_{x}(E) \to E\) is continuous and open, since E is locally path-connected (III, p. 264, corollary). The map \(\varepsilon_{\mathrm{E}}: \lambda_{x}(E) \to E\) that it defines by passing to arc-connected components is therefore continuous and open. The map \(\varepsilon_{\mathrm{E}} \circ (q_{*})^{-1}: \lambda_{b}(B) \to E\) is bijective, continuous, and open. This proves that the topological space \(\lambda_{b}(B)\) , endowed with the topology of compact convergence and the endpoint map, is a covering of B.

(iii)⇒(v). Indeed, the set \(\pi_{1}(B,b)\) , endowed with the quotient topology of the topology of compact convergence, is identified with the fiber of the B-space \(\lambda_{b}(B)\) over b, at the class of the constant loop at b. If \(\lambda_{b}(B)\) is a covering of B, then \(\pi_{1}(B,b)\) is discrete for the topology of compact convergence.

Suppose these assertions hold. By the foregoing, the space \(\lambda_b(B)\) is then a covering of B that is simply path-connected, hence simply connected. The pointed space ( \(\lambda_b(B), \varepsilon_b\) ) is consequently a universal covering of (B, \(b\) ) (I, p. 126, corollary of prop. 3) and the covering \(\lambda_b(B)\) is a Galois covering of B (I, p. 120, prop. 1). The canonical homomorphism \(h_{(\lambda_b(B), \varepsilon_b)}: \pi_1(B, b) \to \text{Aut}_B(\lambda_b(B))\) is then an isomorphism (III, p. 306, prop. 5).

Remark 1. --- Let B be a connected loop-disentanglable space and let b be a point of B. It follows from Theorem 1, (iv), that every action of the group \(\pi_{1}(B,b)\) is admissible. Consequently, every \(\pi_{1}(B,b)\) -set is isomorphic to a \(\pi_{1}(B,b)\) -set \(E_{b}\) , where E is a covering of B (III, p. 318, prop. 7). Recall also (III, p. 310, prop. 2) that if E and \(E'\) are coverings of B, the map \(f \mapsto f_{b}\) induces a bijection from the set of morphisms of B-spaces from E into \(E'\) onto the set of morphisms of \(\pi_{1}(B,b)\) -sets from \(E_{b}\) onto \(E'_{b}\) .

In the language of categories, one says that the functor which, to a covering E of B, associates the fiber \(E_b\) is an equivalence of categories from the category of coverings of B to the category of \(\pi_1(B,b)\) -sets.*

COROLLARY 1. --- Let B be a loop-disentanglable topological space and let b be a point of B. Let E be a connected covering of B and let x be a point of the fiber \(E_{b}\) . The following properties are equivalent:

\begin{enumerate}
\def\labelenumi{(\roman{enumi})}
\item
  The group \(\pi_{1}(\mathrm{E},x)\) is trivial.
\item
  The space E is simply connected.
\item
  The pointed space (E, x) is a universal covering of (B, b).
\end{enumerate}

The implication (i)⇒(ii) has already been proved under the sole hypothesis that the space B is connected and locally path-connected (III, p. 309, cor. 2 of prop. 1), and the implication (ii)⇒(iii) without any hypothesis (I, p. 126, corollary).

Let us prove (iii)⇒(i). By Theorem 1 of IV, p. 342, there exist a covering E' of B and a point \(x'\) of the fiber \(E'_{b}\) such that the group \(\pi_1(E', x')\) is reduced to the identity element. Under hypothesis (iii), there exists

a pointed B-morphism \(f \colon (\mathrm{E}, x) \to (\mathrm{E}', x')\) . Denote by \(p \colon \mathrm{E} \to \mathrm{B}\) and \(p' \colon \mathrm{E}' \to \mathrm{B}\) the projections of coverings \(\mathrm{E}\) and \(\mathrm{E}'\) ; we have \(p = p' \circ f\) therefore \(\pi_1(p, x) = \pi_1(p', x') \circ \pi_1(f, x)\) . Since the group \(\pi_1(\mathrm{E}', x')\) is trivial, the injective homomorphism \(\pi_1(p, x)\) has as its image the identity element. This proves that the group \(\pi_1(\mathrm{E}, x)\) is reduced to the identity element.

COROLLARY 2. --- Let B be a loop-disentanglable topological space and let b be a point of B. Let (E, x) be a universal covering of the pointed space (B, b). For every covering E' of B and every point \(x' \in E'_{b}\), the unique B-morphism f: E → E' such that f(x) = x' makes E a covering of E'. The pointed space (E, x) is then a universal covering of (E', x').

By prop. 7 of I, p. 81, the space E, equipped with the map f, is a covering of E'. The last assertion then follows from corollary 1.

Let B be a loop-disentanglable topological space and let \(a\) , \(b\) be two points of B. The quotient topological space \(\varpi_{a,b}(\mathrm{B})\) is homeomorphic to the space \(\pi_1(\mathrm{B},a)\) , equipped with the quotient topology of the compact-convergence topology (III, p. 293, remark 3), hence is discrete by theorem 1 of IV, p. 342. Consequently, the strict homotopy class of every path in B joining \(a\) to \(b\) is an open subset of \(\Lambda_{a,b}(\mathrm{B})\) . More generally:

PROPOSITION 4. --- Let B be a loop-disentanglable topological space. Strict homotopy is an open equivalence relation in the topological space \(\mathcal{C}_{c}(\mathbf{I};\mathbf{B})\) .

Suppose first that the space B is simply path-connected. Let \(\varphi\colon\Lambda(B)\to B\times B\) be the map which, to a path \(c\) in B, associates the pair \((c(0),c(1))\) consisting of its origin and its endpoint. For two paths in B to be strictly homotopic, it is necessary and sufficient that they have the same origin and the same endpoint (III, p. 292, prop. 2). Since the space B is connected and locally path-connected, the map \(\varphi\) is surjective, continuous, and open (III, p. 262, prop. 10). The strict homotopy relation is therefore open (TG, I, p. 32, prop. 3).

In the general case, let E be a simply path-connected covering of B, nonempty if \(B \neq \varnothing\) . This covering is surjective because B is connected (I, p. 74, prop. 4). Let p denote the projection of the B-space E and let \(\widetilde{p} \colon \Lambda(\mathrm{E}) \to \Lambda(\mathrm{B})\) be the map which, to a path c in E, associates the path \(p \circ c\) . Equipped with the map \(\widetilde{p}\) , \(\Lambda(\mathrm{E})\) is a covering of \(\Lambda(\mathrm{B})\) (III, p. 302, cor. 1 of prop. 3). For two paths in B to be strictly homotopic, it is necessary and sufficient that they be the images under the map \(\widetilde{p}\) of two strictly homotopic paths in E (III, p. 302, prop. 4). Let U be an open subset of the space \(\Lambda(B)\) . The set \((\widetilde{p})(U)\) is open in \(\Lambda(E)\) ; by the first part of the proof, the set \(U'\) saturated with respect to \((\widetilde{p})(U)\) for the strict homotopy relation is open in \(\Lambda(E)\) . Since \(\Lambda(E)\) is a covering of \(\Lambda(B)\) the set \(\widetilde{p}(U')\) is open in \(\Lambda(B)\) . This proves the proposition, since \(\widetilde{p}(U')\) is the saturation of U for the strict homotopy relation.

